
\def\PUDcodigo{EME122}

\def\PUDcodacad{04506.40}

\def\PUDdisciplina{Máquinas Térmicas e de Fluxo}

\def\PUDsemestre{S9}

\def\PUDhoras{80}

%NOVOS ---------------------------------------------------
\def\PUDhorasTeoria{80}
\def\PUDhorasPratica{0}

\def\PUDinterECA{ \hyperref[PUD:EME115]{Mecanismos (04506.25)}
    
\hyperref[PUD:EME114]{Termodinâmica (04506.27)}

\hyperref[PUD:EME118]{Mecânica dos Fluidos (04506.34)}
    
\hyperref[PUD:EME119]{Sistemas Mecânicos (04506.35)}

\hyperref[PUD:EME121]{Transferência de Calor (04506.39)}
}

\def\PUDatuaisECA{---}

\def\PUDrecursos{Projetor multimídia; Quadro branco e pincel.}

%----------------------------------------------------------

\def\PUDcreditos{4}

\def\PUDprerequisito{ \hyperref[PUD:EME118]{EME118} \hyperref[PUD:EME121]{EME121} }

\def\PUDpreacad{ \hyperref[PUD:EME118]{Mecânica dos Fluidos (04506.34)}

\hyperref[PUD:EME121]{Transferência de Calor (04506.39)} }

\def\PUDobjetivos{Conhecer o princípio de funcionamento das máquinas térmicas encontradas na indústria; Ser capaz de fazer uso efetivo do princípio de funcionamento das máquinas térmicas na prática da engenharia.}

\def\PUDementa{Máquinas de Fluxo: Turbomáquinas (Bombas e Turbinas); Máquinas Térmicas: Ciclos de potência e refrigeração com mudança de fase, Ciclos de Potência e refrigeração sem mudança de fase.
%Ciclos termodinâmicos; Motores de combustão interna; Trocadores de calor; Caldeiras; Turbinas; Compressores.
}

\def\PUDconteudo{ & Turbomáquinas: Bombas: Introdução e classificação, A bomba centrífuga, Curvas de desempenho de bombas e leis de semelhança, Bombas de fluxo misto e de fluxo axial (a rotação específica), Combinando as características da bomba e do sistema.
%Revisão da Termodinâmica: Calor, Trabalho, Características dos Sistemas, Propriedades dos fluidos termodinâmicos, 1ª Lei da termodinâmica, 2ª Lei da termodinâmica, Ciclos termodinâmicos.
 \\ 
 & Turbomáquinas: Teoria das Turbinas de ação e de reação, Turbinas Eólicas.
 %Motores de combustão interna: Introdução, Definição de motores de combustão interna, Classificação dos MCI, Vantagens e desvantagens, Definições, Princípio de funcionamento dos motores alternativos, Motor Wankel, Motor Quasiturbine, Veículos híbridos, Ciclos de potência (Otto e Diesel), Componentes dos MCI, Combustíveis, A combustão no motor Diesel, Lubrificação do MCI, Refrigeração do MCI. 
\\ 
 & Máquinas Térmicas: Ciclos de Potência com mudança de fase: Introdução aos Ciclos de Potência, Ciclo Rankine, Efeitos da Pressão e da Temperatura no Ciclo Rankine, O Ciclo com Reaquecimento, O Ciclo Regenerativo e Aquecedores de Água de Alimentação, Afastamento dos Ciclos Reais em Relação aos Ciclos Ideais, Cogeração e outras configurações, Introdução aos ciclos de Refrigeração, Ciclo de Refrigeração por Compressão de Vapor, Fluidos de Trabalho para Sistemas de Refrigeração por Compressão de Vapor, Afastamento do Ciclo de Refrigeração Real de Compressão de Vapor em Relação ao Ciclo Ideal, Configurações de Ciclos de Refrigeração, O Ciclo de Refrigeração por Absorção.
 %Trocadores de Calor: Tipos de trocadores de calor, O coeficiente global de troca de calor, Trocadores de calor compactos, Torres de arrefecimento. 
\\ 
 & Máquinas Térmicas: Ciclos de Potência com Fluidos de Trabalho Gasosos: Ciclos Padrão a Ar, O Ciclo Brayton, O Ciclo Simples de Turbina a Gás com Regenerador, Configurações do Ciclo de Turbina a Gás para Centrais de Potência, O Ciclo Padrão a Ar para Propulsão a Jato, O Ciclo Padrão de Refrigeração a Ar, Ciclos de Potência dos Motores, O Ciclo Otto, O Ciclo Diesel, O Ciclo Stirling, Os Ciclos Atkinson e Miller, Ciclos Combinados de Potência e Refrigeração.
 %Caldeiras: Conceitos básicos de vapor, Gerador de vapor, Histórico da utilização industrial do vapor, Principais componentes das caldeiras, Caldeiras flamotubulares, Caldeiras aquatubulares, Caldeiras elétricas, Segurança, Inspeção da caldeira a vapor, Caldeiras de combustível sólido, Caldeiras de combustíveis líquidos, Cuidados especiais. 
%\\ 
 %& Turbinas: Definição, Aplicações, Classificação,, Turbinas a vapor, Turbinas a gás, Turbinas hidráulicas, Turbinas aeronáuticas, Turbinas eólicas. 
%\\ 
 %& Compressores: Tipos de compressores;  Rendimento volumétrico de espaço nocivo;  O efeito da temperatura de evaporação sobre a vazão de refrigerante;  O efeito da temperatura de evaporação sobre a capacidade frigorífica;  O efeito da temperatura de evaporação sobre a potência de compressão;  O efeito da temperatura de condensação sobre a vazão de refrigerante e a capacidade de refrigeração;  O efeito da temperatura de condensação sobre a potência de compressão;  Catálogos de fabricantes;  Sistema de controle de capacidade linha SMC;  Sistema de controle de capacidade linha CMO;  Sistema de arrefecimento;  Resfriamento de óleo;  Sistema de lubrificação
 }

\def\PUDmetodo{Aulas expositórias que podem ser teóricas e/ou práticas, onde as práticas no laboratório serão marcadas no decorrer da disciplina. A metodologia também contemplará o acompanhamento do desenvolvimento de um programa de simulação de uma máquina térmica ou de fluxo (uma diferente para cada aluno) existente no mercado.}

\def\PUDavalia{A avaliação será feita com aplicação de provas teóricas e/ou práticas, além da possibilidade de inclusão de trabalhos e seminários no decorrer da disciplina. O software de simulação da máquina térmica servirá como uma das avaliações.
}

\def\PUDbibbasica{ & \href{www.biblioteca.ifce.edu.br/index.asp?codigo_sophia=65410}{Borgnakke}, C.; Sonntag, R. E.  Fundamentos da Termodinâmica.  São Paulo, SP: Edgard Blücher, 2013.  728p.  ISBN: 9788521207924. 
\\ 
 & \href{biblioteca.ifce.edu.br/index.asp?codigo_sophia=23814}{White}, Frank M. Mecânica dos fluidos. 6. ed. Porto Alegre: AMGH, 2011. 880 p. ISBN 9788563308214.
\\ 
 &  \href{www.biblioteca.ifce.edu.br/index.asp?codigo_sophia=63098}{Souza}, Zulcy de.  Projeto de máquinas de fluxo: tomo III, turbinas hidráulicas com rotores tipo Francis.  vol. 3.  Rio de Janeiro, RJ: Interciência, 2011.  142p.  ISBN: 9788571932807. }

\def\PUDbibcomplementar{ & \href{www.biblioteca.ifce.edu.br/index.asp?codigo_sophia=63097}{Souza}, Zulcy de.  Projeto de máquinas de fluxo: tomo I, base teórica e experimental.  Rio de Janeiro, RJ: Interciência, 2011.  179p.  ISBN: 9788571932586.
%\\
%&  Boyce, M. P.  Gas Turbine Engineering Handbook, 2ª edição, Editora Science \& Technology Books.  ISBN : 0884157326
\\ 
 &  \href{www.biblioteca.ifce.edu.br/index.asp?codigo_sophia=24058}{Moran}, M. J.; Shapiro, H. N.; Muson, B. R.; DeWitt, D. P.  Introdução à engenharia de sistemas térmicos: termodinâmica, mecânica dos fluidos e transferência de calor.  Rio de Janeiro, RJ: LTC, 2005.  604p.  ISBN: 8521614463.
\\ 
 &  \href{www.biblioteca.ifce.edu.br/index.asp?codigo_sophia=61984}{Incropera}, F. P.; Witt, D. P.  Fundamentos de Transferência de Calor e de Massa.  7ª ed.  Rio de Janeiro, RJ: LTC, 2014.  672p.  ISBN: 9788521625049.
\\
 &  \href{www.biblioteca.ifce.edu.br/index.asp?codigo_sophia=64103}{Boyce}, Meherwan P.  Gas turbine engineering handbook.  4º ed.  Amsterdam, HOL: Elsevier, 2012.  956p.  9780123838421.
\\
 &  \href{www.biblioteca.ifce.edu.br/index.asp?codigo_sophia=24460}{Botelho}, Manoel Henrique Campos.  Operação de caldeiras: gerenciamento, controle e manutenção.  São Paulo, SP: Blucher, 2011.  204p.  ISBN: 9788521205883. 
 \\
 & \href{www.biblioteca.ifce.edu.br/index.asp?codigo_sophia=23804}{Silva}, Napoleão Fernandes da.  Compressores alternativos industriais: teoria e prática.  Rio de Janeiro, RJ: Interciência, 2009.  419p.  ISBN: 9788571932159.}

\def\PUDautor{Francisco Frederico dos Santos Matos}

\def\PUDdata{2014-04-23}

\def\PUDrevisao{3}

\def\PUDrevisor{Francisco Frederico dos Santos Matos}

\def\PUDdatarevisao{2019-05-15}

\PUDcorpo
